A great conjunction is a conjunction of Jupiter and Saturn for observers on
Earth. Assume that Jupiter and Saturn have circular orbits in the ecliptic
plane.

The time between successive conjunctions may vary slightly as viewed from the
Earth. However, the average time period of the great conjunctions is the same
as that of an observer at the centre of the Solar system.

\begin{description}
\item[a)] Find the average great conjunction period (in years) and average
  heliocentric angle between two successive great conjunctions (in degrees).
  \hfill[6]
\item[b)] The next great conjunction will be on 21st December 2020 with an
  elongation of $30.3^\circ$ East of the Sun. Suggest the constellation in
  which the conjunction on 21st December 2020 will occur. (Give the IAU Latin
  name or IAU three-letter abbreviation of the constellation, i.e.\ Ursa Major
  or UMa) \hfill[2]
\end{description}

In 1606, Johannes Kepler determined that in some years the great conjunction
can be happen three times in the same year due to the retrograde motions of
the planets. He also determined that such an event happened in the year 7 BC,
which could have been the event commonly known as ``The Star of Bethlehem''.
For the calculations below you may ignore the precession of the axis of the
Earth.

\begin{description}
\item[c)] Suggest the constellation in which the great conjunctions in 7 BC
  occurred. (Give the IAU Latin name or IAU three-letter abbreviation of the
  constellation, i.e.\ Ursa Major or UMa) \hfill[8]
\item[d)] During the second conjunction of the series of three conjunctions
  in 7 BC, suggest the constellation the Sun was in as viewed by the observer
  on Earth. (Give the IAU Latin name or IAU three-letter abbreviation of the
  constellation, i.e.\ Ursa Major or UMa) \hfill[4]
\end{description}
